% Imporditud: tools/import_eksperimendid.py
% Allikas: Eesti füüsikaolümpiaadi eksperimendi ülesannete
%   kogu (Rael Kalda, 2023), ülesanne Ü93, lahendus L93.
\setAuthor{EFO žürii}
\setRound{lõppvoor}
\setYear{2004}
\setNumber{G E1}
\setDifficulty{3}
\setTopic{staatika}
\setVahendid{korrapäratu kujuga plaat, joonlaud, kaaluviht}
\setPunktid{8}
\setAllikas{Eksperimendi kogu Ü93, lahendus lk 73}
\setLahendusseis{toores}

\prob{Plaat}
Määrake plaadi mass.

\hint

\solu
Määrame esmalt plaadi raskuskeskme. Selleks paigutame plaadi laua servale nii,

et see on peaaegu kukkumas. Märgime laua serva plaadile. Pöörame plaati ning

kordame kirjeldatud tegevust. Saadud kahe sirge lõikepunkt ongi plaadi raskus-

keskme asukoht. Nüüd on juba lihtne määrata plaadi mass kangkaalumise põhi-

mõttel lkoormis lplaat mplaat = mkoormis ,

kus lkoormis ja lplaat on vastavad kangi õlad. Kui koormise massi võib lugeda täpseks, siis maksimaalne relatiivne viga

$\delta$ = $\Delta$lkoormis + $\Delta$lplaat . lkoormis lplaat

Plaadi mass M = mplaat (1 $\pm \delta$).
\probend
